\documentclass[aspectratio=169]{beamer}

\usetheme{Madrid}
\usecolortheme{default}
\definecolor{ucbblue}{RGB}{0, 51, 102}

\setbeamercolor{structure}{fg=ucbblue}
\setbeamercolor*{palette primary}{fg=white,bg=ucbblue}
\setbeamercolor*{palette secondary}{fg=white,bg=ucbblue!80}
\setbeamercolor*{palette tertiary}{fg=white,bg=ucbblue!90}
\setbeamercolor{title}{fg=white, bg=ucbblue}
\setbeamercolor{frametitle}{fg=white, bg=ucbblue}

\usepackage[T1]{fontenc}
\usepackage[utf8]{inputenc}
\usepackage{tgpagella}
\usefonttheme{serif} % Soluciona los errores de fuente "Font shape undefined"
\usepackage[spanish, es-noshorthands]{babel}
\usepackage{amsmath, amssymb}
\usepackage[most]{tcolorbox}

\title[Tarjetas de Error]{Tarjetas de Análisis de Errores: Recuperación Analítica}
\author[IMT-221]{Circuitos Electrónicos II}
\institute{Universidad Católica Boliviana ``San Pablo'' — Sede Tarija}
\date{}

\begin{document}

\begin{frame}
    \titlepage
\end{frame}

\begin{frame}{Diagnóstico de Errores en Circuitos}
    \textbf{Instrucciones para la dinámica[cite: 8]:}
    \vspace{0.4cm}
    
    Analice cada fragmento de cálculo como una cadena:
    \begin{equation*}
        \text{modelo} \rightarrow \text{ecuación} \rightarrow \text{álgebra} \rightarrow \text{resultado} \rightarrow \text{interpretación} \text{[cite: 8]}
    \end{equation*}
    
    \vspace{0.4cm}
    Busque el \textbf{primer punto inválido}, no solamente el resultado final incorrecto[cite: 8]. Responda mentalmente antes de la revelación:
    \begin{enumerate}
        \item ¿Cuál es el primer paso inválido?[cite: 8]
        \item ¿Qué tipo de error es (Tag)?[cite: 8]
        \item ¿Cómo se corrige?[cite: 8]
    \end{enumerate}
\end{frame}

\begin{frame}{Error Card 1 — Capacitor[cite: 8]}
    \textbf{Desarrollo del estudiante[cite: 8]:}
    \begin{align*}
        Z_C &= j\omega C \\
        Z_C &= j(6283)(159 \times 10^{-9}) \\
        Z_C &= j0.001\,\Omega
    \end{align*}
    
    \vspace{0.3cm}
    \pause
    \begin{tcolorbox}[colback=red!5, colframe=red!70!black, title=\textbf{Diagnóstico}, width=\linewidth]
        \textbf{Primer Error:} Definición de impedancia incorrecta[cite: 8]. \\
        \textbf{Corrección:} $Z_C = \frac{1}{j\omega C} = -j\frac{1}{\omega C}$[cite: 8]. \\
        \textbf{Tag:} \texttt{[Impedance] / [Model]}[cite: 8].
    \end{tcolorbox}
\end{frame}

\begin{frame}{Error Card 2 — KCL Sign[cite: 8]}
    El estudiante \textbf{define todas las corrientes saliendo} del nodo $V_A$, pero escribe[cite: 8]:
    \begin{equation*}
        \frac{V_A - V_s}{R_1} + \frac{V_A}{Z_C} - \frac{V_A - V_B}{R_2} = 0 \text{[cite: 8]}
    \end{equation*}

    \vspace{0.3cm}
    \pause
    \begin{tcolorbox}[colback=red!5, colframe=red!70!black, title=\textbf{Diagnóstico}, width=\linewidth]
        \textbf{Primer Error:} La tercera corriente fue definida con un signo negativo, asumiendo implícitamente la dirección contraria a su definición inicial[cite: 8]. \\
        \textbf{Corrección:} $ \frac{V_A - V_s}{R_1} + \frac{V_A}{Z_C} + \frac{V_A - V_B}{R_2} = 0 $[cite: 8]. \\
        \textbf{Tag:} \texttt{[Sign] / [KCL/KVL]}[cite: 8].
    \end{tcolorbox}
\end{frame}

\begin{frame}{Error Card 3 — Magnitudes Prematuras[cite: 8]}
    Dado $Z_R = 1000\,\Omega, Z_C = -j1000\,\Omega$, el estudiante escribe[cite: 8]:
    \begin{equation*}
        Z_R + Z_C = |1000| + |-j1000| = 2000\,\Omega \text{[cite: 8]}
    \end{equation*}

    \vspace{0.3cm}
    \pause
    \begin{tcolorbox}[colback=red!5, colframe=red!70!black, title=\textbf{Diagnóstico}, width=\linewidth]
        \textbf{Primer Error:} Sustituir fasores por magnitudes antes de completar la operación compleja[cite: 8]. \\
        \textbf{Corrección:} $Z_R + Z_C = 1000 - j1000 \implies |Z_R + Z_C| = 1414\,\Omega$[cite: 8]. \\
        \textbf{Tag:} \texttt{[Complex arithmetic]}[cite: 8].
    \end{tcolorbox}
\end{frame}

\begin{frame}{Error Card 4 — Polar Addition[cite: 8]}
    \textbf{Desarrollo del estudiante[cite: 8]:}
    \begin{equation*}
        5\angle0^\circ + 5\angle90^\circ = 10\angle90^\circ \text{[cite: 8]}
    \end{equation*}

    \vspace{0.3cm}
    \pause
    \begin{tcolorbox}[colback=red!5, colframe=red!70!black, title=\textbf{Diagnóstico}, width=\linewidth]
        \textbf{Primer Error:} Sumar magnitudes y ángulos directamente en forma polar[cite: 8]. \\
        \textbf{Corrección:} Convertir a rectangular ($5 + j5$), y luego a polar $\implies 7.07\angle45^\circ$[cite: 8]. \\
        \textbf{Tag:} \texttt{[Complex arithmetic]}[cite: 8].
    \end{tcolorbox}
\end{frame}

\begin{frame}{Error Card 5 — Physical Plausibility[cite: 8]}
    Un divisor pasivo sin fuentes dependientes entrega[cite: 8]:
    \begin{equation*}
        |V_o| = 37\,\text{V} \text{[cite: 8]}
    \end{equation*}
    Cuando la única fuente independiente tiene $|V_s| = 10\,\text{V}$[cite: 8]. \\
    El estudiante concluye: \textit{"La calculadora dio 37 V, por tanto es correcto."[cite: 8]}

    \vspace{0.3cm}
    \pause
    \begin{tcolorbox}[colback=red!5, colframe=red!70!black, title=\textbf{Diagnóstico}, width=\linewidth]
        \textbf{Primer Error:} Aceptar un resultado matemático sin comprobar su coherencia física con los límites del sistema[cite: 8]. \\
        \textbf{Corrección:} Activar una auditoría completa (signos, complejos, KVL)[cite: 8]. \\
        \textbf{Tag:} \texttt{[Physical interpretation]}[cite: 8].
    \end{tcolorbox}
\end{frame}

\begin{frame}{Error Card 6 — Current Direction[cite: 8]}
    El estudiante define[cite: 8]:
    \begin{equation*}
        I = \frac{V_A - V_B}{Z} \text{[cite: 8]}
    \end{equation*}
    Y posteriormente utiliza en la siguiente malla[cite: 8]:
    \begin{equation*}
        V_B - V_A = IZ \text{[cite: 8]}
    \end{equation*}

    \vspace{0.3cm}
    \pause
    \begin{tcolorbox}[colback=red!5, colframe=red!70!black, title=\textbf{Diagnóstico}, width=\linewidth]
        \textbf{Primer Error:} Ambas expresiones implican direcciones de referencia opuestas para la misma variable en el mismo circuito[cite: 8]. \\
        \textbf{Corrección:} Mantener consistencia $\implies V_A - V_B = IZ$[cite: 8]. \\
        \textbf{Tag:} \texttt{[Sign] / [Solution organization]}[cite: 8].
    \end{tcolorbox}
\end{frame}

\end{document}
